\section{机器人大脑：VLM 与 VLA 的双系统}

数据闭环解决“模型吃什么”，机器人大脑则解决“模型怎样行动”。这也是本期从产业分析回到技术结构的地方。

当主持人问“大脑怎么做”时，高继扬把机器人大脑拆成两个基础模型：一个是上层的 VLM，用于指令拆解、逻辑思考和任务规划；另一个是动作基础模型 VLA，用 vision 和 language 产生 action，驱动本体执行任务。这是本期最关键的技术结构。

\begin{figure}[H]
\centering
\includegraphics[width=0.96\textwidth]{robot-brain-dual-system.png}
\caption{机器人大脑的双系统结构：VLM 拆解指令，VLA 产生动作。}
\end{figure}

\begin{knowledgebox}{读图：VLM 与 VLA 为什么要拆开}
VLM 适合处理模糊指令、逻辑拆解和更通用的任务规划；VLA 必须低延迟地产生动作，很多时候要跑在端侧。工业场景中任务集合有限，可能直接调用 VLA 的语言接口；家庭等通用场景更需要 VLM 参与复杂拆解。
\end{knowledgebox}

\subsection{端侧延迟与云端推理}

高继扬特别指出，执行动作的模型如果放在云上，延迟会成为很难解决的问题。因此 VLA 作为动作模型，往往需要部署在端侧。VLM 是否必须实时参与，则取决于场景复杂度。二三十个固定动作的工商业场景，可能不需要每一步都调用 VLM；更通用的家庭场景，VLM 才是不可或缺的组成部分。

\begin{warningbox}{不要把“机器人大脑”想成一个大模型}
机器人大脑不是一个单体 LLM。它至少包含任务拆解、动作生成、端侧控制、感知输入、反馈回流和安全约束。把它简化成“接一个大语言模型”，会低估部署和延迟问题。
\end{warningbox}

\subsection{创业公司相对大厂的优势}

上一节拆开了 VLM 和 VLA，这一节回答竞争问题：大厂有基础模型、算力和人才，创业公司凭什么做机器人大脑？高继扬的答案不是泛泛说“更快”，而是回到数据入口。

高继扬把做好 VLA 的成功要素拆成数据、算法、算力/基础设施、人才。大厂在算力、基础设施和人才上强，但常常缺真实操作数据；有整机能力的创业公司则可能在数据入口上更强。尤其在中国，供应链和硬件迭代能力让创业公司更容易把整机、场景和数据采集打通。

\subsection{本章小结}

机器人大脑的关键不是单纯“谁的大模型更强”，而是能否形成端侧可执行的 VLA、合适的 VLM 上层拆解、真实数据闭环和整机部署能力。

